\documentclass[10pt]{article}
\input{../../notes-style.tex}

\usepackage{multicol}
\usepackage{array}

\title{Ch.~7 Cheat Sheet\\\large Heaps and Priority Queues}
\author{}
\date{}

\begin{document}
\maketitle
\vspace{-2em}

\begin{multicols}{2}

\subsection*{What a heap is}
A \textbf{complete} binary tree with the heap property:
\begin{itemize}
    \item \textbf{Max-heap}: every parent $\ge$ both children. Root is max.
    \item \textbf{Min-heap}: every parent $\le$ both children. Root is min.
\end{itemize}
\textit{Not} a BST: left and right siblings are unordered w.r.t.\
each other. No in-order traversal, no search, no range query.

\subsection*{Array representation}
Zero-indexed, root at 0. For node at index $i$:
\[
\text{parent}(i) = \left\lfloor \tfrac{i-1}{2} \right\rfloor,\
\text{left}(i) = 2i+1,\
\text{right}(i) = 2i+2
\]
Works only because the tree is complete (no gaps).

\subsection*{Why array beats pointers}
\begin{tabular}{@{}l|l@{}}
\textbf{Pointer tree}       & \textbf{Array heap}\\
\hline
3 ptrs / node               & no per-node overhead\\
cache miss per child hop    & children 1--2 cells ahead\\
\texttt{new}/\texttt{delete}& \texttt{vector} amort.\ $O(1)$
\end{tabular}

\subsection*{Percolate up (insert fixup)}
\begin{lstlisting}[language=C++]
while (i > 0) {
    int p = (i - 1) / 2;
    if (h[i] <= h[p]) break;
    std::swap(h[i], h[p]);
    i = p;
}
\end{lstlisting}

\subsection*{Percolate down (extract fixup)}
\begin{lstlisting}[language=C++]
while (2*i + 1 < n) {
    int L = 2*i + 1, R = 2*i + 2, big = L;
    if (R < n && h[R] > h[L]) big = R;
    if (h[i] >= h[big]) break;
    std::swap(h[i], h[big]);
    i = big;
}
\end{lstlisting}
Watch: swap with the \emph{larger} child, not an arbitrary child.

\subsection*{Insert / extract}
\begin{itemize}
    \item \textbf{Insert}: \texttt{push\_back}, then percolate up.
    \item \textbf{Extract}: save root, move last to root,
        \texttt{pop\_back}, percolate down.
\end{itemize}
Both $O(\log n)$; peek $O(1)$.

\columnbreak

\subsection*{Cost table}
\begin{tabular}{@{}l|lll@{}}
               & peek  & pop          & push\\
\hline
sorted array   & $O(1)$& $O(1)$       & $O(n)$\\
unsorted array & $O(n)$& $O(n)$       & $O(1)$\\
balanced BST   & $O(\log n)$ & $O(\log n)$ & $O(\log n)$\\
\textbf{heap}  & $O(1)$ & $O(\log n)$ & $O(\log n)$
\end{tabular}

\subsection*{Heapify (bottom-up build)}
For $i$ from $\lfloor n/2 \rfloor - 1$ down to $0$, call
\texttt{percolateDown(i)}. Total $O(n)$, not $O(n \log n)$:
most nodes are near the bottom and do almost no work
($\sum n/2^{i+1} \cdot i = O(n)$).

\subsection*{Heapsort}
\begin{enumerate}
    \item Heapify in place -- $O(n)$.
    \item For \texttt{end = n-1} down to 1: swap
    \texttt{a[0]} with \texttt{a[end]}, then
    \texttt{percolateDown(0, size=end)}. -- $O(n \log n)$.
\end{enumerate}
In place, $O(1)$ extra, \textbf{not stable}. Worst case
$O(n \log n)$ -- why introsort falls back to it.

\subsection*{Priority queue ADT}
Operations: \texttt{push(x)}, \texttt{pop()} (returns extremum),
\texttt{peek()}, \texttt{size()}, \texttt{empty()}. Usually
implemented as a heap. Uses: Dijkstra / A* / Prim, event
simulation, top-$k$, process scheduling.

\subsection*{\texttt{std::priority\_queue}}
\begin{lstlisting}[language=C++]
std::priority_queue<int> maxpq;            // max-heap
std::priority_queue<int, std::vector<int>,
    std::greater<int>> minpq;              // min-heap
\end{lstlisting}
\texttt{.top()} returns \textbf{const ref} -- no in-place priority
change. To update priority: pop + push, or use indexed heap.

\subsection*{Treap (sneak preview of balance)}
Each node carries (key, random priority).\ BST order on keys,
heap order on priorities. Insert as BST leaf, then rotate up while
priority $>$ parent's. Rotations preserve BST order.\
Expected $O(\log n)$ all ops; worst-case $O(n)$ (unlucky priorities).

\subsection*{Top gotchas}
\begin{itemize}
    \item Index arithmetic off-by-one: \texttt{(i-1)/2}, not
        \texttt{i/2}. Must use int division.
    \item Swapping with the smaller child in
        \texttt{percolateDown} -- creates a new violation below.
    \item Using a heap where you need ordered traversal or
        search (use a BST / \texttt{std::map}).
    \item Treating \texttt{pop()} on empty as safe -- UB in C++.
\end{itemize}

\end{multicols}
\end{document}
